% =============================================================================
% Grammars as Probes — restructured draft
% Target: Interpretability for Discovery workshop, NeurIPS 2026
%         https://interpretability4discovery.github.io/
% CFP requirements implemented here:
%   - NeurIPS 2026 workshop LaTeX template (neurips_2026.sty, from the kit
%     linked on the CFP page), submission (anonymous, line-numbered) mode
%   - <= 5 pages main text, self-contained (refs + appendices excluded)
%   - Responsible-use statement (missing one is grounds for desk rejection)
%   - Double blind: no author names/affiliations, third-person citations
% Content reconstructed from the 2026-08 results draft (NueurIPS2026.pdf);
% \TODO marks = items only the full source has (appendix figures/tables).
% =============================================================================
\documentclass{article}

% Double-blind workshop submission mode (anonymized, line-numbered).
% For camera-ready switch to [dblblindworkshop,final].
% The style loads natbib itself; do not load it again.
\usepackage[dblblindworkshop]{neurips_2026}
\workshoptitle{Interpretability for Discovery}

\usepackage[T1]{fontenc}
\usepackage[utf8]{inputenc}
\usepackage{microtype}
\usepackage{amsmath,amssymb}
\usepackage{booktabs}
\usepackage{multirow}
\usepackage{graphicx}
\usepackage{xcolor}
\usepackage{enumitem}
\usepackage{caption}
\usepackage{url}

% Draft machinery: red TODOs for material to port from the full source.
% Flip to \draftfalse for a clean read; ALL must be resolved before submission.
\newif\ifdraft\drafttrue
\newcommand{\TODO}[1]{\ifdraft\textcolor{red!70!black}{[TODO: #1]}\fi}

\title{Grammars as Probes: Solver-Certified Knowledge\\
Extraction and Localization in Language Models}

\author{%
  Anonymous Author(s)\\
  Affiliation\\
  \texttt{email}\\
}

\begin{document}
\maketitle

\begin{abstract}
Interpretability can turn what a model encodes into knowledge that experts can
test---but only if the extracted knowledge is stated precisely, verified
against data, and shown to be what the model actually uses. We build a
pipeline that meets all three requirements in a controlled domain with a
built-in answer key: Rosetta Stone linguistic puzzles, where a grammar either
regenerates the given translation pairs or does not. A symbolic solver (Answer
Set Programming) derives and \emph{certifies} the minimal grammar rules behind
each puzzle, and each certified rule then becomes a causal probe target inside
\texttt{Llama-3.1-8B} as it solves the puzzle in context. Across four
questions in Gilbertese, Ndebele, and Tshiluba (27 certified rules), we find
that rule implementations stratify by depth---morphosyntactic markers localize
in the earliest layers while semantic lemma selections resolve in late,
output-adjacent layers; that each rule is carried by a compact linear subspace
recovered with subspace Distributed Alignment Search; and that sequential
interventions on independent rules commute (31/31 rule pairs), consistent with
additive, non-interfering composition. Standard last-token probing misses all
of this, producing near-zero causal effect. Because every probed claim is a
machine-verified rule rather than an experimenter-chosen concept, the same
recipe applies where the answer key is missing: certify what the model's
behavior commits it to, then test whether its internals compute it.
\end{abstract}

\section{Introduction}
\label{sec:intro}

Models that match or exceed human experts appear to have internalized
regularities that humans have not articulated, and interpretability offers a
way to read those regularities back out. But an extracted statement of model
knowledge is only useful for discovery if it clears two bars. First,
\emph{verification}: the statement must be checkable against data, not an
artifact of the probing choices---in most interpretability work the
experimenter picks the concept to probe (gender, sentiment, syntax) by
intuition, with no independent guarantee that the chosen concept is the one
the task requires. Second, \emph{internal validity}: the model must actually
compute with the extracted knowledge; a statement that merely fits the model's
outputs may be a post-hoc rationalization of a different
mechanism~\citep{lanham2023,turpin2023}.

We develop and validate a pipeline that clears both bars in a domain where the
ground truth is known, so that the pipeline itself can be audited end to end.
Rosetta Stone linguistic puzzles---the core format of the UK Linguistics
Olympiad (UKLO)~\citep{bozhanov2013,lingoly}---present a handful of bilingual
sentence pairs in an unfamiliar low-resource language and ask for translations
of held-out sentences. They come with built-in verification: a proposed
grammar either regenerates every context pair exactly, or it does not. We use
a symbolic solver (Answer Set Programming, ASP~\citep{clingo}) to derive and
\emph{certify} the minimal grammar rules that solve each puzzle
(\S\ref{sec:grounding}). Each certified rule---an object-agreement suffix, a
tense particle, a noun-class concord---is a discrete, machine-verified,
expert-legible claim, and we take exactly these rules (no others) as causal
probe targets inside \texttt{Llama-3.1-8B} while it solves the puzzle in
context (\S\ref{sec:methods}).

On four questions across three languages (27 certified rules), we show:
\begin{enumerate}[leftmargin=1.6em,itemsep=1pt,topsep=2pt]
\item Dense activation-patching sweeps localize every certified rule, and the
implementations stratify by depth: morphosyntactic agreement markers and tense
particles localize sharply in early layers (L0--L4), whereas semantic lemma
choices consolidate in late decision layers (L20--L31)
(\S\ref{sec:results}).
\item Subspace Distributed Alignment Search (SubDAS) isolates each rule in a
compact ($k\le 8$) linear subspace, and linear subspaces match or beat
nonlinear surrogates (\S\ref{sec:ablations}).
\item Sequential interventions on independent grammatical rules commute
(order-KL $\le 0.0019$ nats; 31/31 rule pairs), consistent with additive
composition without representational cross-talk (\S\ref{sec:results}).
\end{enumerate}
A negative result frames the methodology: probing or patching only the final
prompt token---the field's default recipe---yields near-zero causal effect on
these tasks, because rule features live at earlier morpheme positions
(\S\ref{sec:results}).

For this workshop's agenda, the domain is the point: UKLO grammars are
\emph{known} linguistics, which is what lets us audit every step---the
solver's certificate against the official solutions, and the localization
claims against intervention outcomes. What transfers to domains without an
answer key is the recipe: certify what the model's behavior commits it to (a
program that provably reproduces the observed data), then use the certified
components as probe targets to test which of them the model genuinely
computes. That combination---external verification plus internal causal
validation---is what makes extracted model knowledge a candidate for
scientific discovery rather than a plausible story.

\section{Related work}
\label{sec:related}

\textbf{Linguistic olympiad puzzles as benchmarks.}
LINGOLY~\citep{lingoly} established that frontier LLMs struggle with
olympiad-level linguistic reasoning; LINGOLY-TOO~\citep{lingolytoo} showed via
orthographic obfuscation that part of measured performance is memorization;
related datasets and analyses include modeLing~\citep{modeling} and studies of
where models break down~\citep{unveiling,llmsolvegen}. These report accuracy,
not mechanism; we ask \emph{how} a model that succeeds does so.

\textbf{Program induction and hypothesis search.}
Ellis et al.~\citep{ellis2022} synthesized phonological theories with Bayesian
program induction; LLM-proposed, programmatically verified hypotheses help on
inductive-reasoning tasks~\citep{wang2023hypothesis}. We adopt the
solver-verified program but flip its purpose: the certified grammar is not
(only) a way to answer the puzzle, it is the ground-truth specification of the
knowledge whose internal representation we then test.
Faithful-by-construction pipelines~\citep{lyu2023,logiclm} guarantee the
emitted chain entails the answer; they do not ask whether it reflects the
model's internal computation, which is our focus.

\textbf{Causal interpretability.}
Our tools build on activation patching and causal
tracing~\citep{meng2022rome}, circuit analysis via path
patching~\citep{wang2022ioi,goldowskydill2023}, and Distributed Alignment
Search (DAS)~\citep{geiger2023das,wu2023boundless}, which finds distributed
linear subspaces implementing causal variables. Prior applications choose the
causal variables by hand or from an intuited high-level algorithm; our
variables come from a solver's certificate, removing the experimenter from the
loop. Trace-level chain-of-thought faithfulness
probes~\citep{lanham2023,turpin2023} test whether verbalized reasoning is
load-bearing; our rule-level formulation tests each verified component of the
task's actual solution.

\section{Solver-certified rule grounding}
\label{sec:grounding}

In interpretability practice, what to probe is usually left to human
intuition, with no independent proof that the chosen concepts are the ones the
task requires. We remove this degree of freedom with a symbolic solver. A
grammar is defined in a small language-agnostic schema---a lexicon
(form\,$\leftrightarrow$\,meaning), concatenative morphology and agreement
dependencies, and constituent-ordering constraints (Appendix~\ref{app:schema}).
Given a puzzle's translated context pairs, compiled to ASP facts, the solver
(clingo~\citep{clingo}) searches for the smallest rule set that regenerates
every context pair exactly and uniquely translates the held-out question; a
rule contradicted by even one pair is rejected. A frontier LLM generates the
initial hypothesis space (candidate morpheme segmentations, lexicon entries,
and ordering constraints) directly from the context pairs%
\TODO{verify the exact frozen model identifier for hypothesis-space
generation; the previous draft said ``Claude 3.5 Opus'', which is not a
released model name}, but acceptance is the solver's alone: when it succeeds,
it outputs a certified, minimal set of discrete rules that provably explain
the puzzle data.

Certification turns each rule into a probe target with an answer key. If the
solver certifies that a language marks object number with a verb suffix, we
can construct minimal pairs isolating exactly that rule and ask where
\texttt{Llama-3.1-8B} computes it. Across our four questions the solver
certifies $N=27$ minimal rules: 6 for Gilbertese Q3, 7 for Ndebele Q1, and 8
each for Tshiluba Q7 and Q8 (full schema, certification formulation, and rule
list in Appendix~\ref{app:schema}).

\section{Localization and circuit analysis}
\label{sec:methods}

\textbf{Activation patching sweep.}
For each certified rule $R$ we build a counterfactual minimal pair: a base
prompt $\mathbf{x}_{\text{base}}$ and source prompt $\mathbf{x}_{\text{source}}$
differing exclusively in the feature $R$ governs (all pairs listed in
Appendix~\ref{app:pairs}). Let $\mathbf{h}^{(\ell,p)}_{\text{base}},
\mathbf{h}^{(\ell,p)}_{\text{source}} \in \mathbb{R}^{d}$ be residual-stream
activations at layer $\ell \in \{0,\dots,31\}$ and token position $p$. With
$\Delta = z(\text{token}_{\text{base}}) - z(\text{token}_{\text{source}})$ the
logit difference at the first divergent answer token, we patch
$\mathbf{h}^{(\ell,p)}_{\text{base}} \leftarrow
\mathbf{h}^{(\ell,p)}_{\text{source}}$ and report the normalized causal
indirect effect
\begin{equation}
\mathrm{CIE}_{\mathrm{norm}}(\ell,p) =
\frac{\Delta_{\text{patched}}(\ell,p) - \Delta_{\text{base}}}
     {\Delta_{\text{source}} - \Delta_{\text{base}}},
\end{equation}
where $1.0$ is complete causal recovery of the source feature and $0.0$ no
effect. From each layer$\times$position grid we keep the top spatially
separated causal sites (selection thresholds in Appendix~\ref{app:details}).

\textbf{SubDAS subspace isolation.}
Patching the full 4096-dimensional residual stream can move unrelated
information. To isolate the rule-specific signal we use a subspace variant of
DAS~\citep{geiger2023das,wu2023boundless}: instead of a full $4096 \times
4096$ rotation, SubDAS learns a rank-$k$ orthonormal projection $W$
($k \in \{1,2,4,8\}$; at most $32k$ parameters), swaps only the projection of
the counterfactual difference into that subspace, and leaves all other
directions untouched. $W$ is trained so that the low-rank swap reproduces the
counterfactual target logits (objective and training details in
Appendix~\ref{app:details}).

\textbf{Path patching.}
Given each rule's subspace and location, we map how attention heads read and
route these signals with path patching~\citep{wang2022ioi,goldowskydill2023},
tracing the certified-rule representations from early extraction sites to late
decision layers and keeping a circuit edge only if its normalized effect
clears the threshold in Appendix~\ref{app:details}.

\textbf{Compositionality.}
If the model handles rules independently, intervening on rule $A$ then $B$
should equal $B$ then $A$, and sequential application should recover the joint
intervention $A{+}B$. For each rule pair we construct a four-prompt set (base,
source $A$, source $B$, joint $A{+}B$), apply the two interventions in both
orders, and measure order sensitivity by the KL divergence between the
resulting output distributions, plus agreement of the resulting answers. We
also score the pruned circuit's minimum description length (MDL) to test
whether a small subcircuit carries the computation
(Appendix~\ref{app:details}).

\section{Experimental setup}
\label{sec:setup}

We evaluate four questions across three languages: UKLO 2023 Gilbertese
Q3,\footnote{\url{https://www.uklo.org/wp-content/uploads/2023/03/2023_R1_3-Gilbertese.pdf}}
UKLO 2019 Ndebele
Q1,\footnote{\url{https://www.uklo.org/wp-content/uploads/2022/05/2019_9-Ndebele.pdf}}
and UKLO 2017 Tshiluba Q7 and
Q8\footnote{\url{https://www.uklo.org/wp-content/uploads/2022/05/2017_4.-Tshiluba.pdf}}
(selection rationale and linguistic profiles in Appendix~\ref{app:puzzles}).
We study only questions \texttt{Llama-3.1-8B} answers exactly correctly under
in-context learning, so interventions probe a mechanism that in fact succeeds;
\S\ref{sec:limitations} discusses the selection effect this induces.

All experiments use \texttt{meta-llama/Llama-3.1-8B}~\citep{llama3} (32
layers, $d=4096$, 32 heads), with interventions executed remotely via
NNsight/NDIF~\citep{nnsight}. SubDAS bases are trained with Adam
($\eta=10^{-3}$, 100 steps, ridge $\alpha=1.0$); ranks $k \in \{1,2,4,8\}$ for
Gilbertese Q3 ($N=18$ counterfactual training pairs) and $k \in \{1,2\}$ for
Ndebele Q1 and Tshiluba Q7/Q8 ($N=2$ pairs; $k \le N$). In total the study
comprises 238{,}975 patched forward evaluations, 145 SubDAS-trained causal
sites, 5{,}069 path-patching traces, and 32 dual-intervention runs.

\section{Results}
\label{sec:results}

\textbf{The default recipe fails: last-token probing shows nothing.}
Probing and patching only at the final prompt token with a single direction
---the standard reading-out recipe---produces near-zero causal effect on every
certified rule (Appendix~\ref{app:lasttoken}). The rule features are computed
at earlier morpheme positions rather than pooled at the prompt-final token.
All positive results below required the full layer$\times$position sweep; we
flag this as a caution for knowledge-extraction pipelines built on last-token
probes.

\textbf{Every certified rule localizes, and depth stratifies by rule type.}
Across all 238{,}975 grid evaluations, every solver-certified rule exhibits
localized causal hotspots (full table in Appendix~\ref{app:sweep}). A
consistent pattern holds across all three languages: grammatical concords,
tense markers, and polarity particles localize sharply in the earliest layers,
while semantic lemma selections (noun roots, verb stems) consolidate in late,
output-adjacent layers. Table~\ref{tab:localization} shows representative
peaks.

\begin{table}[t]
\centering
\small
\caption{\textbf{Representative localization peaks} (full sweep results in
Appendix~\ref{app:sweep}). Morphosyntactic rules peak in layers L0--L4;
lemma-selection rules peak in layers L29--L31. Sites are (layer, token
position); $|\mathrm{CIE}_{\mathrm{norm}}|$ can exceed 1 when patching
overshoots the source logit difference.}
\label{tab:localization}
\begin{tabular}{@{}llllr@{}}
\toprule
Question & Rule & Rule type & Peak site & $|\mathrm{CIE}_{\mathrm{norm}}|$ \\
\midrule
Gilbertese Q3 & \texttt{tense}       & morphosyntactic & L00, P001 & 3.39 \\
Ndebele Q1    & \texttt{yes}         & morphosyntactic & L03, P003 & 1.20 \\
Tshiluba Q7   & \texttt{tense}       & morphosyntactic & L00, P253 & 24.00 \\
\addlinespace
Gilbertese Q3 & \texttt{obj\_lemma}  & lemma selection & L30, P195 & 0.98 \\
Ndebele Q1    & \texttt{verb}        & lemma selection & L29, P299 & 1.00 \\
Tshiluba Q8   & \texttt{subj\_lemma} & lemma selection & L31, P261 & 0.99 \\
\bottomrule
\end{tabular}
\end{table}

\begin{figure}[t]
\centering
\fbox{\parbox{0.9\linewidth}{\centering\vspace{2.2cm}
\TODO{port Figure 9 from the full source: layer$\times$position
$\mathrm{CIE}_{\mathrm{norm}}$ heatmaps, (a) early morphosyntactic vs.\
(b) late lemma-selection rules}
\vspace{2.2cm}}}
\caption{\textbf{Layer$\times$position causal maps for certified rules.}
Morphosyntactic rules (a) peak in the earliest layers at morpheme positions;
lemma selections (b) peak in late, output-adjacent layers.}
\label{fig:heatmaps}
\end{figure}

\textbf{Rules live in compact linear subspaces.}
SubDAS shows each rule's causal effect is concentrated in a low-rank subspace
(rank-wise validation in Appendix~\ref{app:subdas}). In Gilbertese Q3, rank
$k=8$ recovers $\mathrm{CIE}_{\mathrm{norm}} = 0.575$ for \texttt{tense}; in
Ndebele Q1 and Tshiluba Q7/Q8, ranks $k \in \{1,2\}$ achieve stable causal
effect ($\mathrm{CIE}_{\mathrm{norm}} \in [0.29, 0.45]$). Bantu concords are
distributed across multi-token prefix-plus-stem sequences, which single-site
interventions only partially capture; intervening at two sites raises recovery
(\S\ref{sec:ablations}).

\textbf{Sparse attention circuits route rule signals.}
Path patching identifies specialized heads carrying rule signals from early
extraction sites to late decision sites (circuit graphs in
Appendix~\ref{app:circuits}): in Gilbertese Q3, head L00.H29 carries the
subject-lemma signal ($\mathrm{CIE}=3.0$) and L11.H12 mediates verb agreement
($\mathrm{CIE}=0.49$); in Ndebele Q1, L00.H23 routes the polarity marker and
L01.H26 subject concord; in Tshiluba Q7/Q8, L04.H16 routes tense prefixes and
L22.H12 diminutive concords. Pruning to these circuits retains
76.5--85.3\% of full-model faithfulness while removing over 85\% of graph
edges (Table~\ref{tab:composition}).

\textbf{Independent rules compose additively.}
Sequential interventions on independent SubDAS operators commute almost
perfectly (Table~\ref{tab:composition}): order-KL divergence is
$D_{\mathrm{KL}} \le 0.0001$ nats across all Ndebele Q1 rule pairs and
$\le 0.0019$ nats across Tshiluba Q7/Q8 canonical pairs, and binary answer
commutativity passes on 31/31 evaluated rule pairs. Sequential application
also recovers the joint intervention. Pruned circuits achieve this at
parsimony ratios below $1.4\times10^{-6}$ of the model's parameter capacity
(Pareto analysis in Appendix~\ref{app:mdl}). Together these results are
consistent with the model composing discrete grammatical rules as
non-interfering, near-orthogonal linear operations on compact subcircuits.

\begin{table}[t]
\centering
\small
\caption{\textbf{Circuit pruning and rule-composition summary} per question.
Faithfulness = pruned-circuit fidelity to the full model; order-KL = maximum
KL divergence between intervention orders across tested rule pairs;
commutativity = binary answer agreement across orders.
\TODO{fill the two missing cells from Tables 4--5 / Appendix H of the full
source; confirm per-question pair counts summing to 31}}
\label{tab:composition}
\begin{tabular}{@{}lrrrr@{}}
\toprule
Question & Edges pruned & Faithfulness & Order-KL (nats) & Commutativity \\
\midrule
Gilbertese Q3 & \TODO{n} & \multirow{4}{*}{\shortstack[r]{76.5--85.3\%\\(range)}} & \TODO{n} & \multirow{4}{*}{\shortstack[r]{31/31\\(100\%)}} \\
Ndebele Q1    & 78  & & $\le 0.0001$ & \\
Tshiluba Q7   & 220 & & $\le 0.0019$ & \\
Tshiluba Q8   & 79  & & $\le 0.0019$ & \\
\bottomrule
\end{tabular}
\end{table}

\section{Ablations}
\label{sec:ablations}

Controls confirming the findings are not artifacts (details in
Appendix~\ref{app:ablations}): (1)~\emph{Linearity:} 2-layer MLP surrogates
trained on the same counterfactual pairs do not beat linear SubDAS, which
matches or outperforms them at 79.2\% of Gilbertese Q3 sites (95/120), 53.5\%
of Ndebele Q1 (38/71), 61.3\% of Tshiluba Q7 (98/160), and 66.2\% of Tshiluba
Q8 (53/80)---added nonlinear capacity does not improve causal steering.
(2)~\emph{Rule specificity:} random orthonormal bases produce negligible
effect ($|\mathrm{CIE}| < 0.01$), and swapping bases between unrelated rules
(e.g., a \texttt{tense} basis at a \texttt{subj\_lemma} site) fails to move
target logits in 80--100\% of runs. (3)~\emph{Global commutativity:} order-KL
measured over the top-50 next-token distribution, not just the answer logit,
stays $\le 0.0001$ nats. (4)~\emph{Multi-site interventions:} Bantu concords
span prefix and stem positions; patching two sites simultaneously raises
Tshiluba Q7 recovery from $\mathrm{CIE}=0.26$ to $0.48$, confirming the
distributed-but-linear picture.

\section{Limitations}
\label{sec:limitations}

Our study is deep, not broad: four questions, three languages, one 8B model.
The selection of questions the model answers exactly correctly is necessary to
study a succeeding mechanism, but it conditions every finding on success---the
mechanisms behind near-misses may differ. Ndebele and Tshiluba rules have only
$N=2$ counterfactual pairs, bounding SubDAS rank at $k\le2$ and limiting
statistical power; results there are directional. Single-token answer
evaluation sidesteps multi-token autoregressive generation. The grammar schema
covers concatenative morphology and consistent ordering; fusional morphology,
reduplication, and phonological alternation are future work. Finally, our
evidence for additive composition is consistency across 31 rule pairs, not a
proof; wider rule inventories could yet reveal interference.

\section{Conclusion}
\label{sec:conclusion}

We showed that solver-certified grammar rules---machine-verified statements of
exactly the knowledge a Rosetta puzzle requires---can serve as probe targets
that make knowledge extraction from a model's internals honest: on four UKLO
questions, \texttt{Llama-3.1-8B} implements each certified rule in a
depth-stratified, low-rank linear subspace, routes it through sparse attention
circuits, and composes independent rules commutatively. The workshop's
question is what models know that we don't; our contribution is procedural:
when a domain admits an executable hypothesis format and a checker, the
certify-then-localize loop turns "the model seems to know $X$" into a pair of
auditable claims---$X$ provably accounts for the observed data, and the model
causally computes $X$---which is the standard extracted knowledge should meet
before an expert spends time testing it in domains where, unlike here, no
answer key exists.

\section*{Responsible-use statement}
This work analyzes the internal mechanisms of a publicly released language
model on educational linguistics puzzles; it introduces no new model, no
generation capability, and no data about individuals. UKLO puzzle materials
are used under their free-for-educational-use terms, and we credit the puzzle
authors via the official sources cited. The principal societal risk we
foresee is epistemic: knowledge-extraction pipelines can lend unearned
authority to spurious "discoveries" read out of models. Our method is designed
to mitigate exactly this failure mode---every extracted claim must be
certified against data by a solver and causally validated in the model---and
we report the failure of the standard last-token recipe so practitioners do
not over-trust cheaper extraction methods. Interventions that steer model
outputs (SubDAS) are rule-specific, puzzle-scoped, and of no plausible dual
use beyond the interpretability setting studied here.

\begingroup
\small
\bibliographystyle{unsrtnat}
\bibliography{references}
\endgroup

\appendix

% ---------------------------------------------------------------------------
% Appendices A-K below mirror the full source draft's structure. Only
% Appendix A is reconstructed here; the rest must be ported verbatim from the
% full LaTeX source (marked TODO). Reviewers are not required to read these;
% the main text stands alone.
% ---------------------------------------------------------------------------

\section{Formal grammar schema and certified rules}
\label{app:schema}

\subsection{Grammar specification}
A grammar $G = (\mathcal{L}, \mathcal{M}, \mathcal{O})$ is defined in a
language-agnostic schema: \textbf{Lexicon} ($\mathcal{L}$): form-to-meaning
mappings between semantic lemmas and surface stems (Gilbertese
\emph{noori-}\,$\leftrightarrow$\,`see', \emph{kune-}\,$\leftrightarrow$\,`find';
Ndebele \emph{-bona}\,$\leftrightarrow$\,`see',
\emph{-thanda}\,$\leftrightarrow$\,`like').
\textbf{Morphology and agreement} ($\mathcal{M}$): concatenative affixes and
agreement dependencies (Bantu subject concord \emph{u-}/\emph{ba-} by noun
class; diminutive \emph{ka-}/\emph{tu-}; Gilbertese object-number verb
suffixes \emph{-a} [sg] vs.\ \emph{-i} [pl]).
\textbf{Ordering constraints} ($\mathcal{O}$): partial orders over clausal
constituents (VOS in Gilbertese; SVO in Ndebele and Tshiluba).

\subsection{Solver certification}
The solver compiles context pairs to ASP facts and constraints and verifies
derivability:
\begin{equation}
\forall (x_i, y_i) \in \mathcal{D}_{\text{context}},\quad
G \vdash \mathrm{Translate}(x_i) = y_i
\quad\text{and}\quad
G \vdash \mathrm{Decode}(x_{\text{query}}) = y_{\text{target}},
\end{equation}
guaranteeing $G$ regenerates every context pair without contradiction and
decodes the held-out question uniquely.

\subsection{Certified rules ($N=27$)}
\textbf{Gilbertese Q3} (6): \texttt{subj\_lemma} (subject noun root),
\texttt{subj\_num} (subject number agreement, \emph{e-} vs.\ \emph{a-}),
\texttt{tense} (\emph{na} future vs.\ $\varnothing$ past), \texttt{verb}
(verb stem), \texttt{obj\_lemma} (object noun root), \texttt{obj\_num}
(object-number verb suffix \emph{-a} vs.\ \emph{-i}).
\textbf{Ndebele Q1} (7): \texttt{yes} (affirmative particle \emph{yebo}),
\texttt{subj\_agr} (noun-class concord \emph{u-}/\emph{ba-}/\emph{si-}),
\texttt{verb} (verb root), \texttt{inf} (infinitive \emph{uku-}),
\texttt{obj} (object NP/suffix), \texttt{adv} (temporal/manner adverbial),
\texttt{part} (participial/aspectual particle).
\textbf{Tshiluba Q7/Q8} (8 each): \texttt{subj\_lemma}, \texttt{subj\_num},
\texttt{subj\_small} (diminutive \emph{ka-}/\emph{tu-}), \texttt{tense}
(\emph{a-}/\emph{aka-}/\emph{aku-}), \texttt{verb}, \texttt{obj\_lemma},
\texttt{obj\_num}, \texttt{obj\_small}.

\section{SubDAS formulation and training objective}
\label{app:details}
\TODO{port from full source (Appendix B), plus the sweep site-selection
thresholds (top-5 sites, $\ge$15\% normalized effect, $\ge$3 layers and
$\ge$5 positions apart), path-patching edge threshold (15\%), and the
commutativity-KL and MDL scoring definitions (Appendix C).}

\section{Puzzle selection and linguistic profiles}
\label{app:puzzles}
\TODO{port from full source (Appendix D).}

\section{Full patching-sweep results}
\label{app:sweep}
\TODO{port from full source (Table 1, Appendix E).}

\section{SubDAS rank validation}
\label{app:subdas}
\TODO{port from full source (Table 2, Appendix F).}

\section{Ablation details}
\label{app:ablations}
\TODO{port from full source (Table 3, Appendix G).}

\section{Commutativity results}
\label{app:commutativity}
\TODO{port from full source (Tables 4--5, Appendix H).}

\section{Last-token probing results}
\label{app:lasttoken}
\TODO{port from full source (Table 6, Appendix I).}

\section{MDL / Pareto analysis}
\label{app:mdl}
\TODO{port from full source (Figure 11, Appendix J).}

\section{Counterfactual minimal pairs}
\label{app:pairs}
\TODO{port from full source (Appendix K).}

\section{Circuit graphs}
\label{app:circuits}
\TODO{port from full source (Figure 10).}

\end{document}
